\section{Discussion}
\label{sec:discussion}

\textbf{Why might the horizon land so differently on the two optimizers?} We can offer
mechanism-shaped observations, not a demonstration. The two losses consume the oracle's $8$
scores differently: \GRPO\ standardises all eight and trains on all eight, so a horizon that
changes \emph{which} sibling looks best --- and by how much relative to the group spread ---
re-weights every gradient; \PTO\ keeps one best-vs-worst pair, so the horizon matters only when
it changes the argmax or argmin, and the DPO update then treats the pair as categorical. The
behavioural asymmetry in \S\ref{sec:behaviour} (only \GRPO\ converts look-ahead into
more-questions) is consistent with that reading. Second, \PTO's greedy trunk feeds each chosen
candidate back into the growing conversation, so a $\Kz$ \PTO\ already receives a weak form of
trajectory feedback that a $\Kz$ \GRPO\ does not --- which would compress the room the explicit
lever has to add. Both stories fit; neither is isolated by this design. The cut that would decide
--- one optimizer trained with the horizon swapped mid-run --- does not exist in these data.

\textbf{Report levers as interaction cells.} The practical upshot of
\S\S\ref{sec:interaction}--\ref{sec:regime} is methodological. A reward-design lever evaluated on
one optimizer in one regime yields one cell of a factorial, and this paper shows a case where the
adjacent cell has the opposite sign and the same cell in an easier regime has the opposite sign
again. Claims of the form ``technique $X$ improves LLM-judge training'' should state the
optimizer family, the budget regime (\S\ref{sec:cost}: the verdict flips with spend), and the
task regime --- or be read as unscoped.

\textbf{The horizon is a safety lever in both families.} The most portable positive finding is
behavioural: turn-level reward taught \emph{both} optimizers unearned praise, and the five-turn
window suppressed it in both, because it moves the exploit's payoff inside the judged span. That
generalised across the optimizer factor even where the score effect did not --- a reason to treat
reward horizon as a first-class safety knob in judge-trained dialogue systems, independent of
whether it buys score.

\textbf{What would change our minds.} A second training run per arm reversing the endpoint
ordering (now the largest unmeasured variance --- the evaluation draws are replicated); an
intermediate $K$ locating a threshold; a mid-run horizon swap isolating the mechanism; or a human
MI coder disagreeing with the over-praise finding.
